\section{Scope and source pins}\label{scope-and-source-pins}

This calculation completes the specific promise in transcript A1545,
node \texttt{baccc895-0427-42ee-b6f6-4c61752dfde2}, to compute ``the
full τ-chart dual, including bottom fibers created by reversed support
arrows.'' It uses the actual four-point chart, coefficient spaces,
differential, and right adjoint of the retained \emph{Arithmetic
cohomology over the τ-base}, \texttt{NOTE.md}, SHA-256
\texttt{d03e71ad18f187bd89880cb8ca6fc9c5d6f260b3de8e8e5702c27db97e7b63c3}.
The complete transcript segment U0041--U0054 has SHA-256
\texttt{e79b46722d9a2ee95b68ae6512f37a19fd41e4b3ecc4e53583cd357d0442c251}.
The reconstruction convention allowing a nonzero bottom fibre is the
fully read \texttt{tex/support\_diagrams.tex}, equations (D1)--(D8). No
generic-stalk restriction replaces global cohomology here.

\section{1. Original objects, signs, and both index
sets}\label{original-objects-signs-and-both-index-sets}

Retain \(E=\mathbb C\), the even Schwartz space \[
V=\{\phi\in\mathcal S(\mathbb R):\phi(-x)=\phi(x),\ \phi(0)=0,\ \widehat\phi(0)=0\},
\qquad
\widehat\phi(\xi)=\int_{\mathbb R}\phi(x)e^{-2\pi i x\xi}\,dx,
\] and \[
\mathscr B=\{F\in C^\infty(\mathbb R_{>0}):
\sup_{x>0}x^b|(-x\partial_x)^kF(x)|<\infty\quad(b\in\mathbb Z,k\ge0)\}.
\] Write \(\mathcal F:V\to V\) for this Fourier transform; on this even
space \(\mathcal F^{-1}=\mathcal F\). Retain \[
\Theta\phi(x)=\sum_{n\ne0}\phi(nx),\qquad
JF(x)=x^{-1}F(1/x),\qquad J\Theta=\Theta\mathcal F,
\qquad Q=\mathscr B/\Theta V,
\] and denote the quotient by \(q:\mathscr B\to Q\). The source proves
injectivity of \(\Theta\) from its Mellin identity in \(\Re s>1\); no
zero in the critical strip is assumed absent.

The geometric index poset remains \(P=\{+,-,\eta,\sigma\}\), with
\(+<\eta<\sigma\) and \(-<\eta<\sigma\). The support index poset is
separately \[
L=\{\varnothing,+,-,J_0\},\qquad J_0=\{+,-\},
\] ordered by inclusion. The symbol \(J_0\) denotes a support mask;
\(J\) above denotes the original inversion operator. At active mask
\(A\), the coefficient diagram \(\mathcal T_A\) retains precisely the
source legs belonging to \(A\), has overlap \(\mathscr B\), and has zero
vector stalk at \(\sigma\). At the empty mask it is the zero diagram.
The source computes \[
\mathsf A_\tau(F)=R\Gamma(P,F)
\simeq[F_+\oplus F_-\xrightarrow{r_+-r_-}F_\eta],
\quad
K_\tau(W)=[I_\eta(W)\longrightarrow I_+(W)\oplus I_-(W)],
\] with degrees \(0,1\) and \(-1,0\), respectively, and differential
signs \(+,-\). Its proved adjunction gives \[
D_A(W):=R\operatorname{Hom}_{\operatorname{Fun}(P,\operatorname{Vect}_E)}
(\mathcal T_A,K_\tau(W))
\cong R\operatorname{Hom}_E(\mathsf A_\tau(\mathcal T_A),W).
\] Here \(W\) is an ordinary coefficient vector space, and
\(X^\vee=\operatorname{Hom}_E(X,W)\) throughout the algebraic
calculation.

\section{2. The complete complexes and every support
arrow}\label{the-complete-complexes-and-every-support-arrow}

The following are literal representatives, not just their cohomology: \[
\begin{array}{c|c|c|c}
A&D_A^{-1}&D_A^0&\delta_A\\ \hline
\varnothing&0&0&0\\
+&\mathscr B^\vee&V^\vee&\ell\mapsto\ell\Theta\\
-&\mathscr B^\vee&V^\vee&\ell\mapsto-\ell\Theta\mathcal F\\
J_0&\mathscr B^\vee&V^\vee\oplus V^\vee&
\ell\mapsto(\ell\Theta,-\ell\Theta\mathcal F).
\end{array}
\] For \(A\subseteq A'\), the arrow \(D_{A'}\to D_A\) is precomposition
with the original coefficient inclusion. On degree \(-1\) it is the
identity between active masks. On degree zero the arrows \(J_0\to+\) and
\(J_0\to-\) are respectively \[
(a,b)\mapsto a,\qquad(a,b)\mapsto b.
\] Both leaf-to-empty arrows are zero, as is the direct joint-to-empty
arrow. Applying either coordinate projection to \(\delta_{J_0}\ell\)
gives the corresponding displayed leaf differential; the maps therefore
commute with the differentials. Both composites from joint to empty are
zero. These checks prove all naturality and composition identities in
\(\operatorname{Fun}(L^{\mathrm{op}},\operatorname{Ch}(E))\).

\section{3. Cohomology, with an explicit quotient and
inverse}\label{cohomology-with-an-explicit-quotient-and-inverse}

For every active mask, the degree-\(-1\) kernel is \[
\{\ell\in\mathscr B^\vee:\ell\Theta=0\}
\xleftarrow[\sim]{q^\vee}Q^\vee,
\qquad q^\vee(\lambda)=\lambda q.
\] The inverse sends \(\ell\) to \([F]\mapsto\ell(F)\). It is
well-defined precisely because \(\ell\Theta=0\). The Fourier isomorphism
makes the condition on the minus face equivalent to the same equation.
All active transition maps on this kernel are the identity in these
\(Q^\vee\) coordinates.

On either single leg, the degree-zero cohomology is zero. Indeed any
linear functional on \(V\) extends to \(\mathscr B\) through the
injective map \(\Theta\). More concretely, the global retraction
constructed in transcript A1694, fully reproved with explicit fixed
cutoffs and estimates in Appendix A below, is a continuous linear map
\(\Lambda:\mathscr B\to V\) with \(\Lambda\Theta=1\). Then \(a\Lambda\)
lifts \(a\in V^\vee\) on the plus face and \(-b\mathcal F^{-1}\Lambda\)
lifts \(b\in V^\vee\) on the minus face.

At the joint mask define \[
\pi_0:V^\vee\oplus V^\vee\longrightarrow V^\vee,
\qquad\pi_0(a,b)=a\mathcal F+b.
\] This evaluates the functional on the original cycle
\(\psi\mapsto(\mathcal F\psi,\psi)\). The composite
\(\pi_0\delta_{J_0}\) is zero. Conversely, if \(a\mathcal F+b=0\), then
\(b=-a\mathcal F\), and \[
\delta_{J_0}(a\Lambda)=(a,-a\mathcal F)=(a,b).
\] The quotient map is surjective, with representative
\(c\mapsto(0,c)\). Its inverse on the quotient sends \(c\) to the class
of \((0,c)\): subtracting this representative from \((a,b)\) gives \[
(a,b)-(0,a\mathcal F+b)=(a,-a\mathcal F)=\delta_{J_0}(a\Lambda).
\] This proves both inverse identities. Consequently the entire
cohomology diagram is \[
\begin{array}{c|cccc}
A&J_0&+&-&\varnothing\\ \hline
H^{-1}D_A&Q^\vee&Q^\vee&Q^\vee&0\\
H^0D_A&V^\vee&0&0&0.
\end{array}
\] The arrows in degree \(-1\) are identity from joint to either leaf
and zero to the empty mask. Every arrow in degree zero is zero. No
degree-zero joint class is discarded.

In particular \(V^\vee\ne0\) when \(W=E\). A concrete continuous
functional is \(\phi\mapsto\phi(1)\). It is nonzero on \[
\phi_*(x)=(4\pi^2x^4-6\pi x^2)e^{-\pi x^2}\in V.
\] The function is even, Schwartz, and vanishes at zero. Integration by
parts gives
\(\int x^4e^{-\pi x^2}dx=\frac{3}{2\pi}\int x^2e^{-\pi x^2}dx\), proving
its integral is zero; its value at one is \(2\pi(2\pi-3)e^{-\pi}\ne0\).

\section{4. What the reversed bottom actually
contains}\label{what-the-reversed-bottom-actually-contains}

In \(L^{\mathrm{op}}\) the bottom is \(J_0\), the join operation is
original intersection, and the top is the original empty mask. Thus the
dual's bottom cohomology fibres are exactly \(Q^\vee\) in degree \(-1\)
and \(V^\vee\) in degree zero. In particular the latter is nonzero for
\(W=E\). The rule ``every bottom fibre is zero'' is therefore
incompatible with this calculation.

The source's full reconstruction retains these fibres. For either
cohomology diagram \(H^iD\), it is \[
M^i=\coprod_{A\in L^{\mathrm{op}}}H^iD_A,
\] with \[
(A,u)+(B,v)=
(A\cap B,\rho_{A,A\cap B}u+\rho_{B,A\cap B}v),
\quad
r^\bullet(A,u)=(A,ru),
\quad
\tau(A,u)=(J_0,0).
\] The global additive zero is \((J_0,0)\), and \(e(A,u)=(A,0)\).
Associativity transports three amplitudes to \(A\cap B\cap C\); equality
of the two iterated transports proves the two sums agree. The zero law
uses the unique map from the bottom and its zero value. Linearity of
transitions proves distributivity, and the fibrewise vector-space scalar
laws prove the remaining semimodule laws. This is the exact
reconstruction, including nonzero vectors in the fibre of its global
zero under multiplication by \(e\).

If an additional external-absence label is desired while retaining every
original reversed label, it can also be constructed explicitly: adjoin a
new bottom \(\bot_\dagger<L^{\mathrm{op}}\), put zero fibre there, and
retain all old fibres and arrows. The new external zero is
\((\bot_\dagger,0)\). The label map sending \(\bot_\dagger\) and \(J_0\)
both to \(J_0\), and fixing every other label, preserves joins and
bottom. With identity maps on old fibres and the zero map on the new
fibre it reconstructs a surjective split-linear comparison to \(M^i\).
Its only additional identification is between the new external zero and
the zero vector at the old bottom \(J_0\). It does not remove any
nonzero bottom vector. This proves the exact relation between the two
declared conventions; it does not silently impose a zero bottom fibre on
the original reversed diagram.

\section{5. Strict natural splitting fails, but the derived splitting
has a concrete
roof}\label{strict-natural-splitting-fails-but-the-derived-splitting-has-a-concrete-roof}

Let \(C=V^\vee\), let \(Z_C\) be the diagram supported at the bottom
\(J_0\) with value \(C\), and let \(H\) be the degree-\(-1\) cohomology
diagram just calculated. The maps \(q^\vee:H[1]\to D\) and
\(\pi:D\to Z_C[0]\), where \(\pi\) is \(\pi_0\) only at the joint mask
in degree zero, are natural cochain maps.

There is no strict natural cochain section \(s:Z_C[0]\to D\) of \(\pi\)
when \(C\ne0\). At the joint mask write \(s(c)=(a(c),b(c))\). Naturality
to the plus leaf forces \(a(c)=0\), since the source at that leaf is
zero. Naturality to the minus leaf forces \(b(c)=0\). Hence \(\pi s=0\),
contradicting \(\pi s=1_C\). This is exactly a strict-section
obstruction.

It is not a derived nonsplitting statement. Here is the actual surviving
derived map. Write \(b=J_0,x=+,y=-,t=\varnothing\) in the reversed
diamond. For a vertex \(v\), define the projective diagram \[
P_v(C)(w)=\begin{cases}C,&v\le w,\\0,&v\nleq w.\end{cases}
\] The augmented resolution of \(Z_C\) is \[
\mathcal R_C:
P_t(C)\xrightarrow{c\mapsto(c,-c)}
P_x(C)\oplus P_y(C)\xrightarrow{(a,b)\mapsto a+b}
P_b(C)\xrightarrow{\epsilon}Z_C,
\] with degrees \(-2,-1,0\) before the augmentation. At \(b\), the
augmentation is identity. At either leaf the preceding map is identity
and the augmented stalk is zero. At \(t\), the stalk sequence is
\(0\to C\to C^2\to C\to0\), with the displayed anti-diagonal and sum.
This proves exactness stalkwise and on all arrows. Each \(P_v(C)\) is
projective because its Hom functor is \(\operatorname{Hom}_E(C,-_v)\),
an exact functor over a field. Thus \(\epsilon:\mathcal R_C\to Z_C[0]\)
is a quasi-isomorphism.

Define \(f:\mathcal R_C\to D\) as follows. In degree zero it is
determined by \[
f_b^0(c)=(0,c),\qquad f_x^0(c)=0,\qquad f_y^0(c)=c,\qquad f_t^0(c)=0.
\] In degree \(-1\), its map from \(P_x(C)\) is zero, and its map from
\(P_y(C)\) is determined at \(y\) by \[
f_y^{-1}(c)=-c\mathcal F^{-1}\Lambda\in\mathscr B^\vee.
\] At \(t\) it is zero; in degree \(-2\) it is zero. These maps are
natural by their projective-source definitions and the zero target at
\(t\). At \(y\), the chain equation is \[
\delta_y(-c\mathcal F^{-1}\Lambda)
=c\mathcal F^{-1}\Lambda\Theta\mathcal F=c.
\] At \(x\) both sides are zero, and at \(t\) both sides are zero. These
verify every chain equation. Finally \(\pi f=\epsilon\), so \[
Z_C[0]\xleftarrow[\simeq]{\epsilon}\mathcal R_C\xrightarrow{f}D
\] is an explicit right inverse of \(\pi\) in the derived category. The
combined cochain map \[
H[1]\oplus\mathcal R_C\longrightarrow D,
\qquad(h,r)\longmapsto q^\vee h+f(r),
\] induces identity on degree \(-1\) cohomology and the displayed
isomorphism on degree zero cohomology. It is therefore a
quasi-isomorphism. This proves the full derived relationship without
replacing the original support arrows by equalities.

\section{6. The original action and the actual defect of this
roof}\label{the-original-action-and-the-actual-defect-of-this-roof}

For \(a>0\), retain \(U_aF(x)=F(x/a)\). On the source plus and minus
legs the actions are \(U_a\) and \(aU_{1/a}\). On \(\mathscr B\) the
action is \(U_a\). Let \(W=\mathcal L_w\) carry \(a^w\), where the
source supplies \(w=0,1\). Then on the displayed dual complex the action
is \[
\ell\longmapsto a^w\ell U_{1/a},
\qquad
(\alpha,\beta)\longmapsto
(a^w\alpha U_{1/a},a^{w-1}\beta U_a).
\] The source identities \(\Theta U_a=U_a\Theta\) and
\(\mathcal F U_a=aU_{1/a}\mathcal F\) verify directly that these
operators commute with every differential and support arrow. On
cohomology they give \[
H^{-1}:\lambda\longmapsto a^w\lambda\overline U_{1/a},
\qquad
H^0:c\longmapsto a^{w-1}cU_a.
\] For \(w=1\), these are the original weight-one dual action on
\(Q^\vee\) and the explicit action \(c\mapsto cU_a\) on the retained
bottom \(V^\vee\).

The roof above need not be strictly equivariant because \(\Lambda\) is
not asserted equivariant. Its degree-zero component is equivariant. At
the minus generator of degree \(-1\), its exact failure is \[
\begin{aligned}
U_a^{D}f_y^{-1}(c)-f_y^{-1}(U_a^{C}c)
&=a^w c\mathcal F^{-1}
\bigl(U_{1/a}\Lambda-\Lambda U_{1/a}\bigr).
\end{aligned}
\] Here the first \(U_{1/a}\) in parentheses acts on \(V\), and the
second on \(\mathscr B\). The formula follows by substitution and
\(U_a\mathcal F^{-1}=a\mathcal F^{-1}U_{1/a}\). Its precomposition with
\(\Theta\) is zero, using \(\Lambda\Theta=1\) and the theta
intertwining. It consequently factors through the exact injection
\(q^\vee:Q^\vee\to\mathscr B^\vee\), with functional \[
[F]\longmapsto a^w c\mathcal F^{-1}
\bigl(U_{1/a}\Lambda-\Lambda U_{1/a}\bigr)F.
\] This is the computed relation between the original action, the chosen
derived roof, and the retained \(Q^\vee\) fibre. No equivariant
splitting is inferred by erasing this term. In this statement each
projective \(P_v(C)\) carries the displayed coefficient action
\(c\mapsto a^{w-1}cU_a\) at every nonzero stalk, and its restriction
maps are identities; its resolution maps therefore intertwine the
actions.

The derived class of this defect is also computable. Applying
\(\operatorname{Hom}(-,H)\) to the displayed projective resolution gives
the complex \[
\operatorname{Hom}_E(C,Q^\vee)
\xrightarrow{\ell\mapsto(\ell,\ell)}
\operatorname{Hom}_E(C,Q^\vee)^2
\longrightarrow0.
\] The last term is zero because \(H_t=0\). Thus \[
\operatorname{Ext}^1(Z_C,H)
\cong\operatorname{Hom}_E(C,Q^\vee),
\qquad[(u,v)]\longmapsto v-u,
\quad
d\longmapsto[(0,d)].
\] The two maps are inverse: their first composite is identity, and
subtracting \((0,v-u)\) from \((u,v)\) gives the diagonal pair
\((u,u)\). The action defect of the chosen roof has zero plus component
and minus component \(q^\vee d_a\), where \[
d_a(c)([F])=
a^w c\mathcal F^{-1}
(U_{1/a}\Lambda-\Lambda U_{1/a})F.
\] It therefore has the exact Ext class \(d_a\). Equivalently, a
nullhomotopy would be a degree-\(-1\) map from \(P_b(C)\) to \(D\),
represented by the same functional on both leaves; its degree-zero
equation forces that functional to annihilate \(\Theta V\). The leaf
equations require its two copies to equal \((0,q^\vee d_a)\). Such a map
exists exactly when \(d_a=0\). This proves that this particular derived
section intertwines the action exactly when the displayed class
vanishes; it does not assert that no other equivariant section can
exist.

The class is the transpose of the original global theta cocycle by an
exact formula. Put \[
s[F]=(1-\Theta\Lambda)F,\qquad k_b=\Lambda U_b s.
\] Every \(F\) has its original decomposition
\(F=\Theta\Lambda F+s qF\). Substitution, with \(\Lambda\Theta=1\) and
\(U_b\Theta=\Theta U_b\), gives \[
U_b\Lambda-\Lambda U_b=-k_b q.
\] Consequently the same class is \[
\boxed{d_a(c)=-a^w c\mathcal F^{-1}k_{1/a}:Q\longrightarrow W.}
\] The minus sign, Fourier inverse, reciprocal parameter, and
coefficient action are retained. The reversed-support defect and the
source's scaling boundary are thus connected by a proved morphism.

\section{7. Continuity and the finite residue inclusion remain
attached}\label{continuity-and-the-finite-residue-inclusion-remain-attached}

For \(W=E\), all the formulas also hold for continuous duals with their
strong topologies. The maps \(\Theta,\mathcal F,q\) and the source's
explicit \(\Lambda\) are continuous, and \(\Lambda\Theta=1\) splits the
original topological sequence. Thus the plus and minus lifts used above
are continuous; all quotient and inverse maps are continuous on the
displayed finite products. Precomposition by a continuous linear map is
strongly continuous because it takes bounded subsets to bounded subsets.
This verifies the asserted topology of every displayed dual transport.
The derived roof assertion in Section 5 is in the declared algebraic
diagram category; no undeclared derived category of all locally convex
spaces is used.

The explicit inverse pair on strong duals is \[
\mathscr B'_b\longrightarrow V'_b\oplus Q'_b,\quad
\ell\longmapsto(\ell\Theta,\ell s),
\qquad
(\alpha,\lambda)\longmapsto\alpha\Lambda+\lambda q.
\] Its composites are identity by \(\Lambda\Theta=1\), \(qs=1\),
\(\Lambda s=0\), \(q\Theta=0\), and \(\Theta\Lambda+sq=1\). Strong
continuity follows directly from \(p_B(T'\ell)=p_{T(B)}(\ell)\) for a
bounded subset \(B\) and continuous linear \(T\). In particular \(q'\)
identifies \(Q'_b\) with the annihilator of \(\Theta V\), with
continuous inverse \(\ell\mapsto\ell s\). The joint relation subspace is
the closed kernel of \(\pi_0\), and \(c\mapsto[(0,c)]\) is its
continuous quotient inverse.

For the original reflection-stable finite arithmetic packet \(A_Z\),
keep \(g=2\xi\), every actual multiplicity, the complete jet map
\(J_Z:Q\to A_Z\), and the full residue form \[
R_Z(f,h)=\sum_{\rho\in Z}\operatorname{Res}_{s=\rho}
\frac{f^\dagger(s)h(s)}{g(s)}\,ds,
\qquad f^\dagger(s)=\overline{f(1-\bar s)}.
\] The source proves its perfectness, \(J_Z\)'s surjectivity and
continuity, and its covariance with \(\mathcal L_1\). Therefore the map
\[
\overline{A_Z}\longrightarrow Q'_b,
\qquad\bar f\longmapsto([F]\mapsto R_Z(f,J_Z[F]))
\] is the same injective continuous residue map at each active vertex of
\(H^{-1}D\). The active support arrows are identities, so its three
naturality squares commute exactly; at the empty vertex both maps are
zero. Every nilpotent remains in the full \(A_Z\) and residue pairing.
The new bottom \(H^0=V'_b\), the action defect of Section 6, and the
original arithmetic residue inclusion now coexist in one completely
specified reversed support calculation.

For the continuous-jet assertion one can retain an explicit estimate.
Write \(r=\Re\rho\) and choose an integer \(n>|r|\). Differentiating the
Mellin integral and splitting it at one gives, for every \(j\ge0\), \[
|(\mathcal MF)^{(j)}(\rho)|
\le j!\left(
\frac{p_{n,0}(F)}{(n-r)^{j+1}}+
\frac{p_{-n,0}(F)}{(n+r)^{j+1}}\right).
\] The two integrals used here are respectively
\(\int_1^\infty x^{r-n}(\log x)^j\,dx/x=j!/(n-r)^{j+1}\) and
\(\int_0^1x^{r+n}|\log x|^j\,dx/x=j!/(n+r)^{j+1}\). They follow from
substitution \(y=|\log x|\) and \(j\) integrations by parts. Thus every
derivative in the finite jet is continuous, with its actual order
retained. The map descends by the original quotient topology since it
annihilates \(\Theta V\).

The finished result is the explicit dual diagram, its nonzero bottom
degree-zero fibre, the precise strict-section obstruction, the
constructed derived roof that survives that obstruction, and its full
scaling defect. None of these computations proves a theta weight bound,
removes a finite-prime trace term, or replaces \(\mathsf A_\tau\) by
restriction to \(\sigma\).

\section{Appendix A: complete retraction on the original theta
source}\label{appendix-a-complete-retraction-on-the-original-theta-source}

\subsection{A.1. Original spaces and the fixed auxiliary
functions}\label{a.1.-original-spaces-and-the-fixed-auxiliary-functions}

Work over \(\mathbb C\), with Fourier convention \[
\widehat\phi(\xi)=\int_{\mathbb R}\phi(x)e^{-2\pi ix\xi}\,dx.
\] Retain the original spaces, with their indicated seminorm topologies:
\[
V=\{\phi\in\mathcal S(\mathbb R):\phi(-x)=\phi(x),\
\phi(0)=0,\ \widehat\phi(0)=0\},
\] \[
\mathscr B=\{F\in C^\infty(\mathbb R_{>0}):
p_{b,k}(F)=\sup_{x>0}x^b|D^kF(x)|<\infty
\quad(b\in\mathbb Z,\ k\ge0)\},\qquad D=-x\partial_x.
\] Use Schwartz seminorms
\(\sigma_{m,j}(u)=\sup_{x\in\mathbb R}(1+|x|)^m|u^{(j)}(x)|\). The
Fourier transform preserves the even Schwartz space and satisfies
\(\mathcal F^2\phi(x)=\phi(-x)\). Consequently it preserves \(V\), since
\(\widehat\phi(0)=\int\phi\) and \(\int\widehat\phi=\phi(0)\).

Fix the following explicit smooth functions for the entire construction:
\[
h(t)=
\begin{cases}
e^{-1/t},&t>0,\\
0,&t\le0,
\end{cases}
\qquad
\chi(x)=
\frac{h(1/16-x^2)}
{h(1/16-x^2)+h(x^2-1/64)},\qquad \alpha=\beta=\chi.
\] All derivatives of \(h\) at zero vanish: each derivative for \(t>0\)
is a polynomial in \(1/t\) times \(e^{-1/t}\), which tends to zero. The
denominator defining \(\chi\) never vanishes, since its two arguments
cannot both be nonpositive. Thus \(\chi\) is real, even, smooth, lies in
\([0,1]\), equals one for \(|x|\le1/8\), and vanishes for \(|x|\ge1/4\).
In particular \[
\|\chi\|_2^2\le\frac12.
\] These are choices of reconstruction data on the original coordinates.
The spaces, Fourier convention and arithmetic theta map remain exactly
those displayed here.

\subsection{A.2. Theta continuity and recovery on both exterior
regions}\label{a.2.-theta-continuity-and-recovery-on-both-exterior-regions}

The original maps are \[
\Theta\phi(x)=\sum_{n\ne0}\phi(nx)=2\sum_{n\ge1}\phi(nx),
\qquad JF(x)=x^{-1}F(1/x).
\] Poisson summation gives \(\Theta\mathcal F=J\Theta\) on \(V\). For
this identity, periodize the Schwartz function \(t\mapsto\phi(xt)\); its
Fourier coefficients are \(x^{-1}\widehat\phi(n/x)\), obtained by
integrating over all translated unit intervals. Both the periodization
and its Fourier series converge absolutely with derivatives. Evaluation
at zero gives \(\sum_n\phi(nx)=x^{-1}\sum_n\widehat\phi(n/x)\). The two
zero-index terms vanish by the original moments.

The theta map is continuous into \(\mathscr B\). Indeed, for \(x\ge1\),
an integer \(N>1\), and \(s_{N,k}(\phi)=\sup_{t>0}t^N|D^k\phi(t)|\), \[
|D^k\Theta\phi(x)|
\le2\zeta(N)s_{N,k}(\phi)x^{-N}.
\] Each \(s_{N,k}\) is bounded by finitely many Schwartz seminorms.
Choosing \(N\ge b\) bounds the part of \(p_{b,k}\) on \(x\ge1\). The
identities \[
D^kJ=J(1-D)^k,\qquad
p_{b,k}(JF)\le\sum_{j=0}^k\binom{k}{j}p_{1-b,j}(F)
\] follow by differentiating \(x^{-1}F(1/x)\). For \(x\le1\), apply the
preceding theta estimate to \(\widehat\phi\) at \(1/x\); choosing also
\(N\ge1-b\) bounds this part by
\(2\zeta(N)\sum_{j=0}^k\binom{k}{j}s_{N,j}(\widehat\phi)\). This proves
every required bound and also continuity of \(J\).

For \(x\ne0\), define the actual arithmetic inverse series \[
\mathcal UF(x)=\frac12\sum_{n\ge1}\mu(n)F(n|x|).
\] Here \(\mu\) is the ordinary Möbius function; the factor \(1/2\)
inverts the sum over both nonzero integer signs. Put \[
P_j(T)=(-1)^jT(T+1)\cdots(T+j-1),\quad P_0=1,\qquad
C_{N,j}(F)=\sum_{k=0}^j|[T^k]P_j(T)|p_{N,k}(F).
\] The identity \(x^j\partial_x^j=P_j(D)\) follows by induction.
Therefore, for \(x>0\) and integer \(N>1\), \[
|\partial_x^j\mathcal UF(x)|
\le\frac{\zeta(N)}2C_{N,j}(F)x^{-N-j}.                 \tag{1}
\] The same absolute bound holds for negative \(x\). It proves locally
uniform convergence of every differentiated series away from zero.

For \(\phi\in V\), the divisor rearrangement is legitimate: for fixed
\(x>0\), \(\sum_{m,n\ge1}|\phi(mnx)|\) is finite, by a Schwartz bound of
degree \(N>1\) and \(\sum_{m,n}(mn)^{-N}=\zeta(N)^2\). Thus \[
\mathcal U\Theta\phi(x)
=\sum_{k\ge1}\left(\sum_{d\mid k}\mu(d)\right)\phi(kx)
=\phi(x).
\] The divisor sum is one for \(k=1\) and zero otherwise: prime
factorization gives the product \(\prod_{p\mid k}(1-1)\) when \(k>1\).
Evenness gives recovery on negative \(x\), and Poisson summation gives
the second exterior recovery: \[
\boxed{\mathcal U\Theta\phi=\phi,\qquad
\mathcal UJ\Theta\phi=\widehat\phi\quad(x\ne0).}         \tag{2}
\] In particular \(\Theta\) is injective, including at zero by
continuity. No information about zeros of zeta in a critical strip
enters this proof.

\subsection{A.3. Fixed contraction and full Schwartz
estimates}\label{a.3.-fixed-contraction-and-full-schwartz-estimates}

On \(L^2(\mathbb R,dx)\), set \[
A=M_\chi,\qquad B=\mathcal F^{-1}M_\chi\mathcal F,\qquad T=AB.
\] Its integral kernel is \(\chi(x)\check\chi(x-y)\). Plancherel and the
change of variable \(z=x-y\) give \[
\|T\|^2\le\|T\|_{\rm HS}^2
=\|\chi\|_2^2\|\check\chi\|_2^2
=\|\chi\|_2^4\le\frac14.
\] Consequently the operator-norm convergent series supplies the actual
inverse: \[
\boxed{(1-T)^{-1}=\sum_{n=0}^{\infty}T^n,\qquad
\|(1-T)^{-1}\|\le2.}                                 \tag{3}
\] Multiplying its partial sums by \(1-T\) leaves \(1-T^{n+1}\), which
tends to the identity, proving both inverse identities.

Define exterior functions, extended by zero near the origin, by \[
a_F=(1-\chi)\mathcal UF,\qquad
b_F=(1-\chi)\mathcal UJF,\qquad
YF=a_F+\chi\mathcal F^{-1}b_F.
\] Here are explicit bounds verifying all their Schwartz seminorms.
Write \(c_0=1\) and \(c_r=\|\chi^{(r)}\|_\infty\) for \(r\ge1\).
Leibniz's rule and (1), with \(N=m+2\), give \[
\sigma_{m,j}(a_F)
\le\frac{\zeta(m+2)}2\,9^m
\sum_{r=0}^j\binom jr c_r\,8^{2+j-r}C_{m+2,j-r}(F).  \tag{4}
\] Indeed every nonzero differentiated term occurs at \(|x|\ge1/8\),
where \(1+|x|\le9|x|\) and \(|x|^{-2-j+r}\le8^{2+j-r}\). The same bound
holds for \(b_F\) with \(JF\) in place of \(F\); the displayed bounds
for \(J\) express it in finitely many original \(p_{b,k}(F)\).
Smoothness across the plateau endpoints follows from smoothness of the
cutoffs and their identically zero interior factors.

Fourier inversion is continuous on these Schwartz seminorms. Explicitly,
integration by parts gives \[
\sigma_{m,j}(\mathcal F^{-1}u)
\le\sum_{r=0}^m\binom mr(2\pi)^{j-r}
\bigl\|\partial_\xi^r(\xi^ju)\bigr\|_1.
\] Each integral on the right is bounded by a finite sum of Schwartz
seminorms, using Leibniz's rule and
\(\int_{\mathbb R}(1+|\xi|)^{-2}d\xi=2\). Together with (4), this proves
that \(Y:\mathscr B\to\mathcal S^{\rm even}\) is continuous.

The operator \(T\) also maps \(L^2\) continuously to Schwartz space. Put
\[
b_j=\|(2\pi\xi)^j\chi(\xi)\|_2
\le\frac{(\pi/2)^j}{\sqrt2},\qquad
A_{m,j}=(5/4)^m
\sum_{r=0}^j\binom jr\|\chi^{(r)}\|_\infty b_{j-r}.
\] Cauchy--Schwarz in the Fourier integral bounds
\(\|\partial^jBf\|_\infty\) by \(b_j\|f\|_2\). Leibniz's rule and the
support of \(\chi\) therefore give \[
\sigma_{m,j}(Tf)\le A_{m,j}\|f\|_2.
\] Set \(\Phi F=(1-T)^{-1}YF\), initially in \(L^2\). The equation
\(\Phi F=YF+T\Phi F\) proves that it is Schwartz, and supplies the
explicit continuity estimate \[
\sigma_{m,j}(\Phi F)
\le\sigma_{m,j}(YF)+2A_{m,j}\|YF\|_2.                 \tag{5}
\] The final \(L^2\) norm is bounded by \(\sqrt2\,\sigma_{1,0}(YF)\).
Every operator preserves parity; hence \(\Phi F\) is even. Equations (2)
give \[
Y\Theta\phi=(1-A)\phi+A(1-B)\phi=(1-T)\phi,
\qquad \Phi\Theta\phi=\phi.                          \tag{6}
\]

\subsection{A.4. The original moment correction and exact
splitting}\label{a.4.-the-original-moment-correction-and-exact-splitting}

Retain \[
\gamma_0(x)=e^{-\pi x^2},\qquad
\gamma_2(x)=2\pi x^2e^{-\pi x^2}.
\] Their moment columns \((\phi(0),\int_{\mathbb R}\phi)\) are \((1,1)\)
and \((0,1)\). For the integral constants, the square of the Gaussian
integral is the planar polar integral
\(2\pi\int_0^\infty re^{-\pi r^2}dr=1\). Integrating
\(\partial_x(xe^{-\pi x^2})\) then gives
\(2\pi\int x^2e^{-\pi x^2}dx=1\). Thus the continuous projection \[
P_Vu=u-u(0)\gamma_0-
\left(\int_{\mathbb R}u-u(0)\right)\gamma_2
\] has both original moments zero and is identity on \(V\). Continuity
follows from \(|u(0)|\le\sigma_{0,0}(u)\) and
\(|\int u|\le2\sigma_{2,0}(u)\). The required map is therefore \[
\boxed{\Lambda=P_V(1-T)^{-1}Y:\mathscr B\longrightarrow V,
\qquad \Lambda\Theta=1_V.}                           \tag{7}
\] All constants and seminorm dependence have been supplied in (1)--(5).
For the original quotient \(q:\mathscr B\to Q=\mathscr B/\Theta V\), put
\(K=1-\Theta\Lambda\). Equation (7) gives \(K^2=K\),
\(\ker K=\Theta V\), and \(\operatorname{im}K=\ker\Lambda\). The image
\(\Theta V\) is closed because it is the kernel of continuous \(K\). The
continuous map \(K\) descends by the quotient topology to \[
s:Q\to\mathscr B,\qquad s[F]=KF,\qquad
qs=1,\quad\Lambda s=0,\quad\Theta\Lambda+sq=1.
\] These equations prove the continuous inverse pair \[
V\oplus Q\longrightarrow\mathscr B,\quad
(\phi,u)\longmapsto\Theta\phi+su,\qquad
F\longmapsto(\Lambda F,qF).
\]

Finally, retain \(U_aF(x)=F(x/a)\). It is continuous because
\(p_{b,k}(U_aF)=a^bp_{b,k}(F)\); its source action preserves \(V\) and
obeys \(U_a\Theta=\Theta U_a\). It therefore induces \(\overline U_a\)
on \(Q\). Put \(k_a=\Lambda U_as\). The decomposition of \(U_as\) gives
\[
U_as=\Theta k_a+s\overline U_a,\qquad
k_{ab}=U_ak_b+k_a\overline U_b,\qquad k_1=0.
\] For the middle identity, substitute the first identity with parameter
\(b\) into \(\Lambda U_aU_bs\), and use \(\Lambda U_a\Theta=U_a\) and
\(\Lambda U_as=k_a\). Writing \(F=\Theta\Lambda F+s qF\) also proves \[
U_a\Lambda-\Lambda U_a=-k_aq.
\] Consequently the reversed-support dual roof's action defect is
exactly \[
\boxed{d_a(c)=-a^wc\mathcal F^{-1}k_{1/a}.}
\] This connects the completed retraction to the original action on the
full \(Q\), retaining its Fourier inverse, weight factor, reciprocal
parameter and sign. The calculation does not set this cocycle to zero.
